\documentclass[aspectratio=169, 11pt]{beamer}
% ============================================================================
% Preambulo compartido para presentaciones Beamer
% Estadistica 2026-II - Pontificia Universidad Javeriana
% Incluir con: % ============================================================================
% Preambulo compartido para presentaciones Beamer
% Estadistica 2026-II - Pontificia Universidad Javeriana
% Incluir con: % ============================================================================
% Preambulo compartido para presentaciones Beamer
% Estadistica 2026-II - Pontificia Universidad Javeriana
% Incluir con: \input{../preambulo_beamer.tex} despues de \documentclass
% ============================================================================

\usepackage[T1]{fontenc}
\usepackage[utf8]{inputenc}
\usepackage[spanish]{babel}
\usepackage{amsmath, amssymb, amsfonts}
\usepackage{booktabs}
\usepackage{graphicx}
\usepackage{tikz}
\usetikzlibrary{shapes.geometric, positioning, arrows.meta, calc}
\usepackage{pgfplots}
\pgfplotsset{compat=1.18}
\usepackage{listings}
\usepackage{xcolor}
\usepackage{hyperref}
\usepackage{multicol}

% --- Tema Beamer (minimalista) ---
\usetheme{default}
\usecolortheme{default}
\usefonttheme{default}
\setbeamertemplate{navigation symbols}{}
\setbeamertemplate{caption}[numbered]
\setbeamertemplate{footline}{%
  \hspace{1em}{\tiny\url{https://github.com/polanco-jaime/Curso-Estadistica-2026-3}}\hfill\usebeamerfont{footline}\insertframenumber/\inserttotalframenumber\hspace{1em}\vspace{0.5em}%
}
\setbeamertemplate{itemize items}[circle]
\setbeamertemplate{enumerate items}[default]

% --- Colores institucionales (sutiles) ---
\definecolor{javeriana}{RGB}{0, 62, 105}
\definecolor{javerianalight}{RGB}{0, 105, 170}
\definecolor{codigobg}{RGB}{245, 245, 245}
\definecolor{codigofg}{RGB}{40, 40, 40}
\definecolor{comentario}{RGB}{100, 100, 100}
\definecolor{keyword}{RGB}{0, 100, 150}
\definecolor{string}{RGB}{180, 60, 30}

\setbeamercolor{title}{fg=javeriana}
\setbeamercolor{frametitle}{fg=javeriana}
\setbeamercolor{structure}{fg=javeriana}
\setbeamercolor{block title}{fg=white, bg=javeriana!75}
\setbeamercolor{block body}{bg=javeriana!5}
\setbeamercolor{block title alerted}{fg=white, bg=red!70!black}
\setbeamercolor{block body alerted}{bg=red!5}
\setbeamercolor{block title example}{fg=white, bg=green!40!black}
\setbeamercolor{block body example}{bg=green!5}
\setbeamercolor{item}{fg=javeriana}

\setbeamerfont{title}{series=\bfseries, size=\Large}
\setbeamerfont{frametitle}{series=\bfseries}

% --- Configuracion de listings para R ---
\lstset{
  language=R,
  basicstyle=\ttfamily\footnotesize,
  backgroundcolor=\color{codigobg},
  commentstyle=\color{comentario}\itshape,
  keywordstyle=\color{keyword}\bfseries,
  stringstyle=\color{string},
  numbers=none,
  frame=single,
  rulecolor=\color{javeriana!30},
  breaklines=true,
  showstringspaces=false,
  columns=flexible,
  morekeywords={library, ggplot, aes, geom_histogram, geom_boxplot,
    geom_point, geom_smooth, geom_density, geom_bar, geom_line,
    geom_vline, geom_hline, geom_errorbar, labs, theme_minimal,
    facet_wrap, scale_fill_manual, scale_color_manual,
    summarise, group_by, filter, mutate, select, arrange,
    read.csv, head, str, summary, mean, median, sd, var, cor,
    lm, t.test, chisq.test, wilcox.test, kruskal.test, aov,
    pnorm, qnorm, dnorm, rnorm, qt, pt, qchisq, pchisq, qf, pf,
    confint, coeftest, vcovHC, bptest, vif, cooks.distance,
    TukeyHSD, table, prop.table}
}

% --- Comandos personalizados ---
\newcommand{\R}{\texttt{R}}
\newcommand{\code}[1]{\texttt{\footnotesize #1}}
\newcommand{\variable}[1]{\texttt{#1}}
\newcommand{\paquete}[1]{\texttt{#1}}

% --- Informacion del curso ---
\institute{Pontificia Universidad Javeriana\\
  Facultad de Ciencias Econ\'omicas y Administrativas}
\date{2026-II}
 despues de \documentclass
% ============================================================================

\usepackage[T1]{fontenc}
\usepackage[utf8]{inputenc}
\usepackage[spanish]{babel}
\usepackage{amsmath, amssymb, amsfonts}
\usepackage{booktabs}
\usepackage{graphicx}
\usepackage{tikz}
\usetikzlibrary{shapes.geometric, positioning, arrows.meta, calc}
\usepackage{pgfplots}
\pgfplotsset{compat=1.18}
\usepackage{listings}
\usepackage{xcolor}
\usepackage{hyperref}
\usepackage{multicol}

% --- Tema Beamer (minimalista) ---
\usetheme{default}
\usecolortheme{default}
\usefonttheme{default}
\setbeamertemplate{navigation symbols}{}
\setbeamertemplate{caption}[numbered]
\setbeamertemplate{footline}{%
  \hspace{1em}{\tiny\url{https://github.com/polanco-jaime/Curso-Estadistica-2026-3}}\hfill\usebeamerfont{footline}\insertframenumber/\inserttotalframenumber\hspace{1em}\vspace{0.5em}%
}
\setbeamertemplate{itemize items}[circle]
\setbeamertemplate{enumerate items}[default]

% --- Colores institucionales (sutiles) ---
\definecolor{javeriana}{RGB}{0, 62, 105}
\definecolor{javerianalight}{RGB}{0, 105, 170}
\definecolor{codigobg}{RGB}{245, 245, 245}
\definecolor{codigofg}{RGB}{40, 40, 40}
\definecolor{comentario}{RGB}{100, 100, 100}
\definecolor{keyword}{RGB}{0, 100, 150}
\definecolor{string}{RGB}{180, 60, 30}

\setbeamercolor{title}{fg=javeriana}
\setbeamercolor{frametitle}{fg=javeriana}
\setbeamercolor{structure}{fg=javeriana}
\setbeamercolor{block title}{fg=white, bg=javeriana!75}
\setbeamercolor{block body}{bg=javeriana!5}
\setbeamercolor{block title alerted}{fg=white, bg=red!70!black}
\setbeamercolor{block body alerted}{bg=red!5}
\setbeamercolor{block title example}{fg=white, bg=green!40!black}
\setbeamercolor{block body example}{bg=green!5}
\setbeamercolor{item}{fg=javeriana}

\setbeamerfont{title}{series=\bfseries, size=\Large}
\setbeamerfont{frametitle}{series=\bfseries}

% --- Configuracion de listings para R ---
\lstset{
  language=R,
  basicstyle=\ttfamily\footnotesize,
  backgroundcolor=\color{codigobg},
  commentstyle=\color{comentario}\itshape,
  keywordstyle=\color{keyword}\bfseries,
  stringstyle=\color{string},
  numbers=none,
  frame=single,
  rulecolor=\color{javeriana!30},
  breaklines=true,
  showstringspaces=false,
  columns=flexible,
  morekeywords={library, ggplot, aes, geom_histogram, geom_boxplot,
    geom_point, geom_smooth, geom_density, geom_bar, geom_line,
    geom_vline, geom_hline, geom_errorbar, labs, theme_minimal,
    facet_wrap, scale_fill_manual, scale_color_manual,
    summarise, group_by, filter, mutate, select, arrange,
    read.csv, head, str, summary, mean, median, sd, var, cor,
    lm, t.test, chisq.test, wilcox.test, kruskal.test, aov,
    pnorm, qnorm, dnorm, rnorm, qt, pt, qchisq, pchisq, qf, pf,
    confint, coeftest, vcovHC, bptest, vif, cooks.distance,
    TukeyHSD, table, prop.table}
}

% --- Comandos personalizados ---
\newcommand{\R}{\texttt{R}}
\newcommand{\code}[1]{\texttt{\footnotesize #1}}
\newcommand{\variable}[1]{\texttt{#1}}
\newcommand{\paquete}[1]{\texttt{#1}}

% --- Informacion del curso ---
\institute{Pontificia Universidad Javeriana\\
  Facultad de Ciencias Econ\'omicas y Administrativas}
\date{2026-II}
 despues de \documentclass
% ============================================================================

\usepackage[T1]{fontenc}
\usepackage[utf8]{inputenc}
\usepackage[spanish]{babel}
\usepackage{amsmath, amssymb, amsfonts}
\usepackage{booktabs}
\usepackage{graphicx}
\usepackage{tikz}
\usetikzlibrary{shapes.geometric, positioning, arrows.meta, calc}
\usepackage{pgfplots}
\pgfplotsset{compat=1.18}
\usepackage{listings}
\usepackage{xcolor}
\usepackage{hyperref}
\usepackage{multicol}

% --- Tema Beamer (minimalista) ---
\usetheme{default}
\usecolortheme{default}
\usefonttheme{default}
\setbeamertemplate{navigation symbols}{}
\setbeamertemplate{caption}[numbered]
\setbeamertemplate{footline}{%
  \hspace{1em}{\tiny\url{https://github.com/polanco-jaime/Curso-Estadistica-2026-3}}\hfill\usebeamerfont{footline}\insertframenumber/\inserttotalframenumber\hspace{1em}\vspace{0.5em}%
}
\setbeamertemplate{itemize items}[circle]
\setbeamertemplate{enumerate items}[default]

% --- Colores institucionales (sutiles) ---
\definecolor{javeriana}{RGB}{0, 62, 105}
\definecolor{javerianalight}{RGB}{0, 105, 170}
\definecolor{codigobg}{RGB}{245, 245, 245}
\definecolor{codigofg}{RGB}{40, 40, 40}
\definecolor{comentario}{RGB}{100, 100, 100}
\definecolor{keyword}{RGB}{0, 100, 150}
\definecolor{string}{RGB}{180, 60, 30}

\setbeamercolor{title}{fg=javeriana}
\setbeamercolor{frametitle}{fg=javeriana}
\setbeamercolor{structure}{fg=javeriana}
\setbeamercolor{block title}{fg=white, bg=javeriana!75}
\setbeamercolor{block body}{bg=javeriana!5}
\setbeamercolor{block title alerted}{fg=white, bg=red!70!black}
\setbeamercolor{block body alerted}{bg=red!5}
\setbeamercolor{block title example}{fg=white, bg=green!40!black}
\setbeamercolor{block body example}{bg=green!5}
\setbeamercolor{item}{fg=javeriana}

\setbeamerfont{title}{series=\bfseries, size=\Large}
\setbeamerfont{frametitle}{series=\bfseries}

% --- Configuracion de listings para R ---
\lstset{
  language=R,
  basicstyle=\ttfamily\footnotesize,
  backgroundcolor=\color{codigobg},
  commentstyle=\color{comentario}\itshape,
  keywordstyle=\color{keyword}\bfseries,
  stringstyle=\color{string},
  numbers=none,
  frame=single,
  rulecolor=\color{javeriana!30},
  breaklines=true,
  showstringspaces=false,
  columns=flexible,
  morekeywords={library, ggplot, aes, geom_histogram, geom_boxplot,
    geom_point, geom_smooth, geom_density, geom_bar, geom_line,
    geom_vline, geom_hline, geom_errorbar, labs, theme_minimal,
    facet_wrap, scale_fill_manual, scale_color_manual,
    summarise, group_by, filter, mutate, select, arrange,
    read.csv, head, str, summary, mean, median, sd, var, cor,
    lm, t.test, chisq.test, wilcox.test, kruskal.test, aov,
    pnorm, qnorm, dnorm, rnorm, qt, pt, qchisq, pchisq, qf, pf,
    confint, coeftest, vcovHC, bptest, vif, cooks.distance,
    TukeyHSD, table, prop.table}
}

% --- Comandos personalizados ---
\newcommand{\R}{\texttt{R}}
\newcommand{\code}[1]{\texttt{\footnotesize #1}}
\newcommand{\variable}[1]{\texttt{#1}}
\newcommand{\paquete}[1]{\texttt{#1}}

% --- Informacion del curso ---
\institute{Pontificia Universidad Javeriana\\
  Facultad de Ciencias Econ\'omicas y Administrativas}
\date{2026-II}

\title{Sesión 8: Chi-Cuadrado: Independencia y Bondad de Ajuste}
\subtitle{Probando asociación entre variables categóricas}
\author{Prof.\ Jaime Polanco Jim\'enez}
\date{Vie 16 oct 2026}

\begin{document}

\begin{frame}
\titlepage
\end{frame}

\begin{frame}{Agenda}
\tableofcontents
\end{frame}

% ============================================================================
\section{Tablas de contingencia}
% ============================================================================

\begin{frame}{¿Qué es una tabla de contingencia?}
\begin{block}{Definición}
Una \textbf{tabla de contingencia} (o tabla cruzada) resume la distribución conjunta de dos variables categóricas, mostrando las frecuencias observadas de cada combinación de categorías.
\end{block}

\vspace{0.3cm}
\textbf{Estructura básica:}
\begin{itemize}
  \item Filas: categorías de la variable 1
  \item Columnas: categorías de la variable 2
  \item Celdas: número de observaciones en cada combinación
  \item Márgenes: totales por fila y columna
\end{itemize}

\vspace{0.3cm}
\begin{exampleblock}{Ejemplo en nuestro contexto}
¿Cómo se relaciona el nivel de inglés MCER (A-, A1, A2, B1, B+) con el tipo de IES (pública/privada) entre estudiantes de Negocios Internacionales?
\end{exampleblock}
\end{frame}

\begin{frame}[fragile]{Construyendo una tabla de contingencia en R}
\begin{lstlisting}
# Cargar datos ICFES para Negocios Internacionales
icfes_ni <- read.csv("datos/icfes_ni_2025.csv")

# Crear tabla de contingencia: nivel MCER x tipo IES
tabla_obs <- table(icfes_ni$NIVEL_MCER, icfes_ni$TIPO_IES)
tabla_obs

#             privada publica
#   A-            245     312
#   A1            580     698
#   A2            892     756
#   B1            623     384
#   B+            410     185
\end{lstlisting}

\begin{alertblock}{Interpretación}
Hay 892 estudiantes de IES privadas con nivel A2 en inglés.
\end{alertblock}
\end{frame}

\begin{frame}[fragile]{Frecuencias marginales y relativas}
\textbf{Frecuencias marginales} (totales por fila/columna):
\begin{lstlisting}
# Totales por fila (nivel MCER)
margin.table(tabla_obs, 1)
#   A-   A1   A2   B1   B+
#  557 1278 1648 1007  595

# Totales por columna (tipo IES)
margin.table(tabla_obs, 2)
# privada publica
#    2750    2335
\end{lstlisting}

\vspace{0.2cm}
\textbf{Frecuencias relativas}:
\begin{lstlisting}
# Proporciones del total
prop.table(tabla_obs)

# Proporciones por fila (% dentro de cada nivel MCER)
prop.table(tabla_obs, 1)

# Proporciones por columna (% dentro de cada tipo IES)
prop.table(tabla_obs, 2)
\end{lstlisting}
\end{frame}

\begin{frame}[fragile]{Visualización: mosaico}
\begin{lstlisting}
# Grafico de mosaico
mosaicplot(tabla_obs,
           main = "Nivel MCER x Tipo IES",
           xlab = "Nivel MCER",
           ylab = "Tipo IES",
           color = c("steelblue", "coral"))

# Alternativa: barplot apilado
barplot(tabla_obs,
        legend = TRUE,
        col = c("#1f77b4", "#ff7f0e", "#2ca02c", "#d62728", "#9467bd"),
        main = "Distribucion por tipo de IES",
        xlab = "Tipo IES",
        ylab = "Frecuencia")
\end{lstlisting}

\begin{block}{Observación visual}
Parece haber una diferencia: las IES privadas tienen mayor proporción de estudiantes en niveles B1 y B+. ¿Es estadísticamente significativa?
\end{block}
\end{frame}

% ============================================================================
\section{Prueba $\chi^2$ de independencia}
% ============================================================================

\begin{frame}{Hipótesis de independencia}
\begin{block}{Pregunta de investigación}
¿El nivel de inglés MCER es \textbf{independiente} del tipo de IES?
\end{block}

\vspace{0.3cm}
\textbf{Hipótesis:}
\begin{itemize}
  \item $H_0$: El nivel MCER y el tipo de IES son \textbf{independientes}\\
        (no hay asociación entre las dos variables)
  \item $H_1$: Existe \textbf{asociación} entre el nivel MCER y el tipo de IES\\
        (las variables están relacionadas)
\end{itemize}

\vspace{0.3cm}
\begin{alertblock}{Nivel de significancia}
Usaremos $\alpha = 0.05$ (estándar en ciencias sociales).
\end{alertblock}
\end{frame}

\begin{frame}{Frecuencias esperadas bajo independencia}
Si las dos variables fueran \textbf{independientes}, ¿cuántas observaciones esperaríamos en cada celda?

\vspace{0.2cm}
\begin{block}{Fórmula}
La frecuencia esperada en la celda $(i,j)$ es:
\[
E_{ij} = \frac{R_i \times C_j}{n}
\]
donde:
\begin{itemize}
  \item $R_i$ = total de la fila $i$
  \item $C_j$ = total de la columna $j$
  \item $n$ = total general
\end{itemize}
\end{block}

\vspace{0.2cm}
\textbf{Ejemplo:} celda (A2, privada)
\[
E_{\text{A2,privada}} = \frac{1648 \times 2750}{5085} = \frac{4532000}{5085} \approx 891.24
\]

Observamos 892, esperábamos 891.24 (muy cercano).
\end{frame}

\begin{frame}{El estadístico $\chi^2$}
Medimos la \textbf{discrepancia total} entre frecuencias observadas y esperadas:

\begin{block}{Estadístico de prueba}
\[
\chi^2 = \sum_{i=1}^{r} \sum_{j=1}^{c} \frac{(O_{ij} - E_{ij})^2}{E_{ij}}
\]
donde:
\begin{itemize}
  \item $O_{ij}$ = frecuencia observada en celda $(i,j)$
  \item $E_{ij}$ = frecuencia esperada en celda $(i,j)$
  \item $r$ = número de filas
  \item $c$ = número de columnas
\end{itemize}
\end{block}

\vspace{0.2cm}
\textbf{Interpretación:}
\begin{itemize}
  \item $\chi^2$ cercano a 0 $\Rightarrow$ observado $\approx$ esperado $\Rightarrow$ independencia plausible
  \item $\chi^2$ grande $\Rightarrow$ gran discrepancia $\Rightarrow$ evidencia contra independencia
\end{itemize}
\end{frame}

\begin{frame}{Distribución $\chi^2$ y grados de libertad}
Bajo $H_0$ (independencia), el estadístico $\chi^2$ sigue aproximadamente una distribución chi-cuadrado con:

\[
\text{gl} = (r - 1)(c - 1)
\]

\vspace{0.2cm}
\textbf{En nuestro ejemplo:}
\begin{itemize}
  \item $r = 5$ niveles MCER (A-, A1, A2, B1, B+)
  \item $c = 2$ tipos IES (privada, pública)
  \item gl $= (5-1)(2-1) = 4 \times 1 = 4$
\end{itemize}

\vspace{0.3cm}
\begin{block}{Regla de decisión}
Rechazamos $H_0$ si:
\[
\chi^2_{\text{calc}} > \chi^2_{\alpha, \text{gl}}
\]
o equivalentemente, si el p-valor $< \alpha$.
\end{block}
\end{frame}

\begin{frame}{Condición de validez}
\begin{alertblock}{Requisito importante}
La aproximación $\chi^2$ es válida si:
\[
E_{ij} \geq 5 \quad \text{en al menos el 80\% de las celdas}
\]
\end{alertblock}

\vspace{0.2cm}
Si esta condición NO se cumple:
\begin{itemize}
  \item Combinar categorías con frecuencias bajas
  \item Usar la prueba exacta de Fisher (solo para tablas $2 \times 2$)
  \item Aumentar el tamaño muestral
\end{itemize}

\vspace{0.3cm}
\textbf{En nuestro ejemplo:}\\
Todas las celdas tienen $E_{ij} > 100$, así que la condición se cumple ampliamente.
\end{frame}

\begin{frame}[fragile]{Aplicación en R: \texttt{chisq.test()}}
\begin{lstlisting}
# Realizar la prueba chi-cuadrado
resultado <- chisq.test(tabla_obs)
resultado

#   Pearson's Chi-squared test
#
# data:  tabla_obs
# X-squared = 186.42, df = 4, p-value < 2.2e-16

# Ver las frecuencias esperadas
resultado$expected

#             privada   publica
#   A-      301.20     255.80
#   A1      691.06     586.94
#   A2      891.24     756.76
#   B1      544.40     462.60
#   B+      321.60     272.90
\end{lstlisting}
\end{frame}

\begin{frame}{Interpretación del resultado}
\begin{block}{Output de R}
\begin{itemize}
  \item $\chi^2 = 186.42$
  \item gl $= 4$
  \item p-valor $< 2.2 \times 10^{-16}$ (prácticamente 0)
\end{itemize}
\end{block}

\vspace{0.3cm}
\textbf{Conclusión:}
\begin{itemize}
  \item Como p-valor $< 0.05$, rechazamos $H_0$.
  \item Existe evidencia \textbf{muy fuerte} de que el nivel de inglés MCER y el tipo de IES \textbf{NO son independientes}.
  \item Hay una asociación estadísticamente significativa entre ambas variables.
\end{itemize}

\vspace{0.3cm}
\begin{alertblock}{Importante}
Rechazar $H_0$ indica \textbf{asociación}, pero no implica causalidad.\\
La prueba NO nos dice qué tan \textbf{fuerte} es la asociación.
\end{alertblock}
\end{frame}

\begin{frame}[fragile]{Visualizando las diferencias}
\begin{lstlisting}
# Calcular residuos estandarizados
residuos <- resultado$residuals
residuos

#             privada   publica
#   A-        -3.24      3.51
#   A1        -4.23      4.58
#   A2         0.03     -0.03
#   B1         3.37     -3.65
#   B+         4.93     -5.34

# Los residuos estandarizados > |2| indican celdas
# con diferencias importantes respecto a independencia
\end{lstlisting}

\vspace{0.2cm}
\textbf{Interpretación:}
\begin{itemize}
  \item IES privadas tienen \textbf{más} estudiantes B1 y B+ de lo esperado (+3.37, +4.93)
  \item IES privadas tienen \textbf{menos} estudiantes A- y A1 de lo esperado (-3.24, -4.23)
  \item La situación es simétrica para IES públicas
\end{itemize}
\end{frame}

\begin{frame}[fragile]{Visualización gráfica de residuos}
\begin{lstlisting}
# Mapa de calor de residuos estandarizados
library(ggplot2)
library(reshape2)

residuos_df <- melt(resultado$residuals)
names(residuos_df) <- c("Nivel", "TipoIES", "Residuo")

ggplot(residuos_df, aes(x = TipoIES, y = Nivel, fill = Residuo)) +
  geom_tile(color = "white") +
  scale_fill_gradient2(low = "blue", mid = "white", high = "red",
                       midpoint = 0, limits = c(-6, 6)) +
  geom_text(aes(label = round(Residuo, 2)), size = 5) +
  labs(title = "Residuos estandarizados: Nivel MCER x Tipo IES",
       x = "Tipo de IES", y = "Nivel MCER") +
  theme_minimal()
\end{lstlisting}

\begin{block}{Lectura del gráfico}
Azul = menos de lo esperado; Rojo = más de lo esperado
\end{block}
\end{frame}

% ============================================================================
\section{V de Cramér}
% ============================================================================

\begin{frame}{Midiendo la fuerza de la asociación}
La prueba $\chi^2$ nos dice \textbf{si hay asociación}, pero no \textbf{qué tan fuerte} es.

\vspace{0.2cm}
\begin{block}{V de Cramér}
Es una medida de asociación para tablas de contingencia:
\[
V = \sqrt{\frac{\chi^2}{n \times (\min(r,c) - 1)}}
\]
donde:
\begin{itemize}
  \item $\chi^2$ = estadístico chi-cuadrado
  \item $n$ = tamaño muestral total
  \item $r$ = número de filas
  \item $c$ = número de columnas
\end{itemize}
\end{block}

\vspace{0.2cm}
\textbf{Rango:} $V \in [0, 1]$
\begin{itemize}
  \item $V = 0$ $\Rightarrow$ independencia completa
  \item $V = 1$ $\Rightarrow$ asociación perfecta
\end{itemize}
\end{frame}

\begin{frame}[fragile]{Cálculo de V de Cramér en R}
\begin{lstlisting}
# Calcular V de Cramer
chi2 <- resultado$statistic
n <- sum(tabla_obs)
min_dim <- min(nrow(tabla_obs), ncol(tabla_obs))

V <- sqrt(chi2 / (n * (min_dim - 1)))
V

# [1] 0.1915

# Usando el paquete DescTools (alternativa)
library(DescTools)
CramerV(tabla_obs)

# [1] 0.1915
\end{lstlisting}

\vspace{0.2cm}
\begin{exampleblock}{Cálculo manual}
\[
V = \sqrt{\frac{186.42}{5085 \times (2-1)}} = \sqrt{\frac{186.42}{5085}} = \sqrt{0.0367} \approx 0.1915
\]
\end{exampleblock}
\end{frame}

\begin{frame}{Interpretación de V de Cramér}
\textbf{Reglas de pulgar generales:}
\begin{itemize}
  \item $V < 0.1$ $\Rightarrow$ asociación \textbf{muy débil}
  \item $0.1 \leq V < 0.3$ $\Rightarrow$ asociación \textbf{moderada}
  \item $V \geq 0.3$ $\Rightarrow$ asociación \textbf{fuerte}
\end{itemize}

\vspace{0.3cm}
\textbf{En nuestro ejemplo:}
\begin{itemize}
  \item $V = 0.1915$ (entre 0.1 y 0.3)
  \item Asociación \textbf{moderada} entre nivel MCER y tipo de IES
  \item Estadísticamente significativa (p-valor $< 0.001$), pero la magnitud del efecto es moderada
\end{itemize}

\vspace{0.3cm}
\begin{alertblock}{Significancia estadística vs práctica}
Una asociación puede ser estadísticamente significativa (p-valor pequeño) pero de magnitud modesta en términos prácticos. Siempre reportar ambos: p-valor Y tamaño del efecto.
\end{alertblock}
\end{frame}

\begin{frame}[fragile]{Ejemplo adicional: Estrato × Región}
\begin{lstlisting}
# El estrato socioeconomico es independiente de la region?
tabla_estrato_region <- table(icfes_ni$ESTRATO, icfes_ni$REGION)

chi_est_reg <- chisq.test(tabla_estrato_region)
chi_est_reg

# X-squared = 523.67, df = 20, p-value < 2.2e-16

# V de Cramer
V_est_reg <- sqrt(chi_est_reg$statistic /
                  (sum(tabla_estrato_region) *
                   (min(dim(tabla_estrato_region)) - 1)))
V_est_reg

# [1] 0.226

# Asociacion moderada y estadisticamente significativa
\end{lstlisting}
\end{frame}

% ============================================================================
\section{Prueba $\chi^2$ de bondad de ajuste}
% ============================================================================

\begin{frame}{¿Qué es bondad de ajuste?}
\begin{block}{Pregunta diferente}
En la prueba de \textbf{independencia}, comparamos dos variables categóricas observadas.

En la prueba de \textbf{bondad de ajuste}, comparamos una variable observada con una distribución teórica esperada.
\end{block}

\vspace{0.3cm}
\textbf{Hipótesis:}
\begin{itemize}
  \item $H_0$: Los datos provienen de la distribución especificada
  \item $H_1$: Los datos NO provienen de esa distribución
\end{itemize}

\vspace{0.3cm}
\begin{exampleblock}{Ejemplos de preguntas}
\begin{itemize}
  \item ¿Los niveles MCER se distribuyen uniformemente?
  \item ¿El puntaje de inglés sigue una distribución normal?
  \item ¿La proporción de estudiantes por región coincide con la distribución poblacional de Colombia?
\end{itemize}
\end{exampleblock}
\end{frame}

\begin{frame}{Estadístico de bondad de ajuste}
\begin{block}{Fórmula}
\[
\chi^2 = \sum_{i=1}^{k} \frac{(O_i - E_i)^2}{E_i}
\]
donde:
\begin{itemize}
  \item $O_i$ = frecuencia observada en la categoría $i$
  \item $E_i$ = frecuencia esperada en la categoría $i$ (según la distribución teórica)
  \item $k$ = número de categorías
\end{itemize}
\end{block}

\vspace{0.2cm}
\textbf{Grados de libertad:}
\[
\text{gl} = k - p - 1
\]
donde $p$ = número de parámetros estimados de los datos.

\begin{itemize}
  \item Si la distribución teórica está completamente especificada: $p = 0$
  \item Si estimamos la media y desviación estándar de los datos: $p = 2$
\end{itemize}
\end{frame}

\begin{frame}[fragile]{Ejemplo 1: Distribución uniforme de niveles MCER}
\textbf{Pregunta:} ¿Los 5 niveles MCER se distribuyen uniformemente?

\begin{lstlisting}
# Frecuencias observadas
niveles <- table(icfes_ni$NIVEL_MCER)
niveles

#   A-   A1   A2   B1   B+
#  557 1278 1648 1007  595

# Bajo uniformidad, esperariamos n/5 en cada nivel
n <- sum(niveles)
esperadas_unif <- rep(n / 5, 5)
esperadas_unif

# [1] 1017 1017 1017 1017 1017

# Prueba chi-cuadrado
chisq.test(niveles, p = rep(1/5, 5))

# X-squared = 574.23, df = 4, p-value < 2.2e-16
\end{lstlisting}

\textbf{Conclusión:} Rechazamos $H_0$. La distribución NO es uniforme.
\end{frame}

\begin{frame}[fragile]{Visualización: observado vs esperado (uniforme)}
\begin{lstlisting}
# Comparar graficamente
library(ggplot2)

df_comp <- data.frame(
  Nivel = names(niveles),
  Observado = as.numeric(niveles),
  Esperado = esperadas_unif
)

df_comp_long <- melt(df_comp, id.vars = "Nivel")

ggplot(df_comp_long, aes(x = Nivel, y = value, fill = variable)) +
  geom_bar(stat = "identity", position = "dodge") +
  labs(title = "Distribucion de niveles MCER: observado vs uniforme",
       x = "Nivel MCER", y = "Frecuencia") +
  scale_fill_manual(values = c("steelblue", "coral")) +
  theme_minimal()
\end{lstlisting}

\begin{block}{Observación}
Hay un pico en A2 y menor frecuencia en los extremos (A- y B+).
\end{block}
\end{frame}

\begin{frame}[fragile]{Ejemplo 2: ¿El puntaje de inglés sigue una distribución normal?}
\textbf{Estrategia:}
\begin{enumerate}
  \item Dividir el rango de \variable{MOD\_INGLES\_PUNT} en intervalos
  \item Calcular frecuencias observadas en cada intervalo
  \item Calcular frecuencias esperadas bajo normalidad con $\mu$ y $\sigma$ muestrales
  \item Aplicar prueba $\chi^2$ de bondad de ajuste
\end{enumerate}

\vspace{0.2cm}
\begin{lstlisting}
# Parametros muestrales
media <- mean(icfes_ni$MOD_INGLES_PUNT, na.rm = TRUE)
desv <- sd(icfes_ni$MOD_INGLES_PUNT, na.rm = TRUE)

media  # 158.3
desv   # 22.7

# Crear intervalos (e.g., 10 intervalos equidistantes)
breaks <- seq(min(icfes_ni$MOD_INGLES_PUNT, na.rm = TRUE),
              max(icfes_ni$MOD_INGLES_PUNT, na.rm = TRUE),
              length.out = 11)
\end{lstlisting}
\end{frame}

\begin{frame}[fragile]{Cálculo de frecuencias esperadas bajo normalidad}
\begin{lstlisting}
# Frecuencias observadas
obs_counts <- table(cut(icfes_ni$MOD_INGLES_PUNT,
                        breaks = breaks,
                        include.lowest = TRUE))

# Probabilidades bajo N(media, desv^2)
probs <- numeric(10)
for(i in 1:10) {
  probs[i] <- pnorm(breaks[i+1], mean = media, sd = desv) -
              pnorm(breaks[i], mean = media, sd = desv)
}

# Frecuencias esperadas
n <- sum(!is.na(icfes_ni$MOD_INGLES_PUNT))
esp_counts <- n * probs

# Prueba chi-cuadrado (df = 10 - 2 - 1 = 7)
chisq.test(obs_counts, p = probs)

# X-squared = 45.82, df = 7, p-value = 2.1e-07
\end{lstlisting}

\textbf{Conclusión:} Rechazamos $H_0$. La distribución del puntaje NO es normal.
\end{frame}

\begin{frame}[fragile]{Alternativa: Test de Shapiro-Wilk}
Para probar normalidad en variables continuas, el test de \textbf{Shapiro-Wilk} es más potente:

\begin{lstlisting}
# Test de Shapiro-Wilk
# (solo funciona con n <= 5000, usar una muestra)
set.seed(123)
muestra <- sample(icfes_ni$MOD_INGLES_PUNT, 5000)

shapiro.test(muestra)

#   Shapiro-Wilk normality test
#
# data:  muestra
# W = 0.9823, p-value < 2.2e-16
\end{lstlisting}

\vspace{0.2cm}
\begin{alertblock}{Interpretación}
\begin{itemize}
  \item $H_0$: los datos provienen de una distribución normal
  \item p-valor $< 0.001$ $\Rightarrow$ rechazamos $H_0$
  \item Los puntajes de inglés NO siguen una distribución normal
\end{itemize}
\end{alertblock}
\end{frame}

\begin{frame}[fragile]{Visualización: Q-Q plot}
\begin{lstlisting}
# Q-Q plot para evaluar normalidad visualmente
qqnorm(icfes_ni$MOD_INGLES_PUNT,
       main = "Q-Q Plot: Puntaje de Ingles",
       pch = 20, col = "steelblue")
qqline(icfes_ni$MOD_INGLES_PUNT, col = "red", lwd = 2)

# Si los datos fueran normales, los puntos deberian
# seguir aproximadamente la linea roja
\end{lstlisting}

\vspace{0.2cm}
\begin{block}{Interpretación del Q-Q plot}
\begin{itemize}
  \item Puntos alineados con la línea $\Rightarrow$ normalidad
  \item Desviaciones en las colas $\Rightarrow$ distribución con colas más pesadas/ligeras
  \item Patrón en S $\Rightarrow$ asimetría
\end{itemize}
\end{block}

En nuestro caso: las colas superiores se desvían, indicando asimetría negativa.
\end{frame}

\begin{frame}[fragile]{Ejemplo 3: Distribución regional esperada}
\textbf{Pregunta:} ¿La distribución regional de estudiantes NI refleja la distribución poblacional de Colombia?

\begin{lstlisting}
# Frecuencias observadas
obs_region <- table(icfes_ni$REGION)

# Proporciones poblacionales de Colombia (ejemplo hipotetico)
# Andina: 34%, Caribe: 21%, Pacifica: 18%, Orinoquia: 3%,
# Amazonia: 2%, Bogota: 22%
prop_poblacion <- c(Andina = 0.34, Caribe = 0.21,
                     Pacifica = 0.18, Orinoquia = 0.03,
                     Amazonia = 0.02, Bogota = 0.22)

# Prueba chi-cuadrado
chisq.test(obs_region, p = prop_poblacion)

# X-squared = 128.56, df = 5, p-value < 2.2e-16
\end{lstlisting}

\textbf{Conclusión:} La distribución regional de estudiantes NI NO refleja la distribución poblacional (posiblemente por concentración de programas en ciertas regiones).
\end{frame}

% ============================================================================
\section{Comparación de múltiples proporciones}
% ============================================================================

\begin{frame}{Extensión: k proporciones}
\textbf{Pregunta:} ¿La proporción de estudiantes que supera B1 es igual en las 5 regiones principales?

\vspace{0.2cm}
Esto es equivalente a una tabla de contingencia $2 \times k$:
\begin{itemize}
  \item Fila 1: supera B1 (B1 o B+)
  \item Fila 2: no supera B1 (A-, A1, A2)
  \item $k$ columnas: una por cada región
\end{itemize}

\vspace{0.2cm}
\begin{block}{Hipótesis}
\begin{itemize}
  \item $H_0$: $p_1 = p_2 = \cdots = p_k$ (todas las proporciones son iguales)
  \item $H_1$: al menos una proporción es diferente
\end{itemize}
\end{block}

\vspace{0.2cm}
Usamos la prueba $\chi^2$ de independencia (equivalente) o \texttt{prop.test()} en R.
\end{frame}

\begin{frame}[fragile]{Aplicación: prop.test() para k proporciones}
\begin{lstlisting}
# Crear variable binaria: supera B1
icfes_ni$SUPERA_B1 <- ifelse(icfes_ni$NIVEL_MCER %in% c("B1", "B+"),
                             1, 0)

# Tabla 2x5: supera B1 x region
tabla_b1_region <- table(icfes_ni$SUPERA_B1, icfes_ni$REGION)

# Prueba chi-cuadrado de independencia
chisq.test(tabla_b1_region)

# X-squared = 38.72, df = 4, p-value = 7.8e-08

# Alternativa: prop.test
prop.test(tabla_b1_region)

# 5-sample test for equality of proportions without continuity correction
# X-squared = 38.72, df = 4, p-value = 7.8e-08
\end{lstlisting}

\textbf{Conclusión:} La proporción que supera B1 difiere significativamente entre regiones.
\end{frame}

\begin{frame}[fragile]{Visualización de proporciones por región}
\begin{lstlisting}
# Calcular proporciones por region
prop_b1_region <- prop.table(tabla_b1_region, 2)

# Grafico de barras
barplot(prop_b1_region,
        main = "Proporcion que supera B1 por region",
        xlab = "Region",
        ylab = "Proporcion",
        legend = c("No supera B1", "Supera B1"),
        col = c("coral", "steelblue"),
        beside = FALSE)
\end{lstlisting}

\vspace{0.2cm}
\begin{exampleblock}{Observación}
Bogota y la region Andina muestran las proporciones mas altas de estudiantes que superan B1.
\end{exampleblock}
\end{frame}

% ============================================================================
\section{Resumen U3 completa}
% ============================================================================

\begin{frame}{Diagrama de flujo: ¿Qué prueba usar?}
\begin{center}
\begin{tikzpicture}[node distance=1.5cm, auto,
                    block/.style={rectangle, draw, fill=blue!10,
                                  text width=5em, align=center,
                                  rounded corners, minimum height=3em},
                    decision/.style={diamond, draw, fill=orange!10,
                                     text width=4em, align=center,
                                     inner sep=0pt}]

\node [decision] (tipo) {Tipo de variable};
\node [block, below left=1cm and 1cm of tipo] (num) {Numérica};
\node [block, below right=1cm and 1cm of tipo] (cat) {Categórica};

\node [decision, below=0.8cm of num] (ngrupos) {¿Cuántos grupos?};
\node [block, below left=0.8cm and 0.5cm of ngrupos] (ttest) {1 o 2:\\ t-test};
\node [block, below right=0.8cm and 0.5cm of ngrupos] (anova) {3+:\\ ANOVA};

\node [decision, below=0.8cm of cat] (pregunta) {¿Qué pregunta?};
\node [block, below left=0.8cm and 0.5cm of pregunta] (indep) {Asociación:\\ $\chi^2$ indep.};
\node [block, below right=0.8cm and 0.5cm of pregunta] (bondad) {Bondad:\\ $\chi^2$ ajuste};

\draw[->] (tipo) -- (num);
\draw[->] (tipo) -- (cat);
\draw[->] (num) -- (ngrupos);
\draw[->] (ngrupos) -- (ttest);
\draw[->] (ngrupos) -- (anova);
\draw[->] (cat) -- (pregunta);
\draw[->] (pregunta) -- (indep);
\draw[->] (pregunta) -- (bondad);

\end{tikzpicture}
\end{center}
\end{frame}

\begin{frame}{Resumen: Pruebas para variables numéricas}
\begin{block}{Comparación de medias}
\begin{itemize}
  \item \textbf{1 media:} t-test de una muestra\\
    $H_0: \mu = \mu_0$ vs $H_1: \mu \neq \mu_0$ (o $>$, $<$)

  \item \textbf{2 medias independientes:} t-test de dos muestras\\
    $H_0: \mu_1 = \mu_2$ vs $H_1: \mu_1 \neq \mu_2$

  \item \textbf{2 medias pareadas:} t-test pareado\\
    $H_0: \mu_d = 0$ vs $H_1: \mu_d \neq 0$

  \item \textbf{3+ medias:} ANOVA de un factor\\
    $H_0: \mu_1 = \mu_2 = \cdots = \mu_k$ vs $H_1:$ al menos una diferente
\end{itemize}
\end{block}

\vspace{0.2cm}
\textbf{Alternativas no paramétricas} (si los supuestos de normalidad no se cumplen):
\begin{itemize}
  \item 2 grupos: Mann-Whitney U (Wilcoxon rank-sum)
  \item 3+ grupos: Kruskal-Wallis
\end{itemize}
\end{frame}

\begin{frame}{Resumen: Pruebas para variables categóricas}
\begin{block}{Prueba $\chi^2$ de independencia}
\begin{itemize}
  \item Compara dos variables categóricas
  \item $H_0$: las variables son independientes
  \item Estadístico: $\chi^2 = \sum \frac{(O_{ij} - E_{ij})^2}{E_{ij}}$
  \item gl $= (r-1)(c-1)$
  \item Tamaño del efecto: V de Cramér
\end{itemize}
\end{block}

\begin{block}{Prueba $\chi^2$ de bondad de ajuste}
\begin{itemize}
  \item Compara una variable categórica con una distribución teórica
  \item $H_0$: los datos siguen la distribución especificada
  \item Estadístico: $\chi^2 = \sum \frac{(O_i - E_i)^2}{E_i}$
  \item gl $= k - p - 1$
  \item Casos especiales: uniformidad, normalidad (con Shapiro-Wilk)
\end{itemize}
\end{block}
\end{frame}

\begin{frame}{Recomendaciones finales}
\begin{alertblock}{Siempre verificar supuestos}
\begin{itemize}
  \item $\chi^2$: $E_{ij} \geq 5$ en al menos 80\% de celdas
  \item t-test: normalidad (o $n$ grande), homocedasticidad
  \item ANOVA: normalidad, homocedasticidad, independencia
\end{itemize}
\end{alertblock}

\begin{block}{Reportar resultados completos}
\begin{itemize}
  \item Estadístico de prueba y grados de libertad
  \item p-valor
  \item Tamaño del efecto (V de Cramér, Cohen's d, $\eta^2$, etc.)
  \item Intervalos de confianza cuando sea aplicable
  \item Interpretación en contexto
\end{itemize}
\end{block}

\begin{exampleblock}{Próxima sesión}
Regresión lineal simple (OLS): modelando la relación entre dos variables numéricas.
\end{exampleblock}
\end{frame}

\end{document}
